\documentclass[pdflatex,sn-mathphys-num]{sn-jnl}

\usepackage{graphicx}
\usepackage{amsmath}
\usepackage{booktabs}
\usepackage{siunitx}
\usepackage{multirow}

\jyear{2026}

\begin{document}

\title[A Calibrated Surrogate Framework for PCS Trend Evaluation]{A Calibrated Surrogate Framework for Probabilistic Constellation Shaping Trend Evaluation in Coherent Optical Link Scenarios}

\author*[1]{\fnm{First} \sur{Author}}\email{author@example.com}
\affil*[1]{\orgdiv{Department}, \orgname{Institution}, \orgaddress{\country{Country}}}

\abstract{
Probabilistic constellation shaping (PCS) is a practical method for improving the information-rate efficiency of coherent optical communication systems, but evaluating PCS trends across multiband link scenarios can require a careful balance between modelling cost, reproducibility and physical fidelity. This paper presents a calibrated surrogate framework for evaluating PCS gain trends in coherent optical link scenarios. The framework combines repeated-seed numerical evaluation, symbol-count stability checks and external GNPy-based calibration diagnostics. The work explicitly separates raw absolute agreement from calibrated diagnostic agreement. In the current evidence set, raw absolute GSNR differs from GNPy by a conservative near-constant offset, while a single offset-calibrated diagnostic aligns the C-band multi-span GSNR trend with sub-dB error. The strongest validated C-band diagnostic reports a calibrated RMSE of 0.284458 dB and leave-one-out RMSE of 0.379277 dB. The manuscript therefore positions the method as a reproducible calibrated surrogate workflow rather than a fully physical Manakov/NLSE split-step Fourier transmission simulator or an experimentally validated system model. The resulting framework is useful for transparent PCS trend evaluation, baseline comparison and reproducibility-oriented system-performance studies, while highlighting the additional physical validation required before making direct absolute-performance claims.
}

\keywords{coherent optical communication, probabilistic constellation shaping, calibrated surrogate modelling, GNPy, GSNR, multiband optical links}

\maketitle

\section{Introduction}

Coherent optical communication systems increasingly rely on advanced modulation, digital signal processing and link-level modelling to evaluate capacity and reach under practical transmission constraints. Probabilistic constellation shaping (PCS) is attractive because it can improve information-rate efficiency by adapting the transmitted symbol distribution. However, comparing PCS trends across link conditions is challenging when the modelling workflow mixes physical approximations, learned mitigation blocks, heuristic noise scaling and external reference tools.

This paper addresses that gap by presenting a reproducible calibrated surrogate framework for PCS trend evaluation. The central goal is not to claim direct physical Manakov/NLSE simulation fidelity. Instead, the goal is to provide a transparent workflow that reports PCS trends, stability behaviour and calibrated agreement against external reference diagnostics.

The main contributions are:
\begin{itemize}
    \item a reproducible calibrated surrogate workflow for PCS trend evaluation in coherent optical link scenarios;
    \item repeated-seed and symbol-count stability checks that reduce the risk of reporting short-sequence or seed-specific PCS artifacts;
    \item explicit separation between raw absolute GSNR disagreement and offset-calibrated external diagnostic agreement;
    \item a GNPy-based C-band external calibration diagnostic showing that the raw GSNR mismatch is largely explained by a conservative near-constant offset across span count;
    \item a baseline structure that separates uniform modulation, PCS-only behaviour, neural-only mitigation where available, combined PCS-plus-neural configurations and external-reference diagnostics;
    \item claim-controlled reporting that states the current physical limitations rather than presenting the surrogate as a direct Manakov/NLSE simulator.
\end{itemize}

\section{Related Work}

PCS has been widely studied as a rate-adaptation and reach-improvement technique for coherent optical transmission systems. Prior work has shown that shaping the transmitted symbol distribution can improve capacity and enable flexible operation without necessarily changing the modulation format \cite{buchali2015pcs64qam,yankov2016constellation}. These studies motivate PCS as a meaningful degree of freedom for link-level trend evaluation.

Analytical nonlinear-propagation models provide another important reference point. The Gaussian-noise (GN) model is widely used to approximate nonlinear interference in uncompensated coherent optical systems \cite{poggiolini2012gn}. The enhanced GN model further refines nonlinear-interference prediction and discusses accuracy--complexity trade-offs \cite{carena2014egn}. These models are useful baselines because they provide physically motivated references for link-quality estimation without requiring exhaustive waveform simulation.

GNPy provides an open-source physical-layer-aware planning and quality-of-transmission workflow for optical networks \cite{ferrari2020gnpy}. In this paper, GNPy is used as an external calibration diagnostic rather than as proof of direct absolute agreement. This distinction is important because the current framework exhibits a conservative raw GSNR offset relative to GNPy before calibration.

Nonlinear compensation has also been studied using digital backpropagation and machine-learning-based processing. Digital backpropagation jointly addresses dispersion and nonlinear impairments in coherent detection systems \cite{ip2008dbp}. Machine-learning methods have been investigated for optical performance monitoring, nonlinear compensation and receiver-side processing \cite{zibar2016ml,lau2018mlopm,tizikara2021opm,bakhshali2023nnnlic}. These works motivate including neural-only and combined PCS-plus-neural configurations as baselines when the corresponding data are available.

The present manuscript differs from these prior directions by focusing on claim-controlled, reproducible calibrated trend evaluation. The objective is not to replace full physical modelling or measurement-based assessment, but to provide a transparent workflow that reports raw mismatch, calibrated diagnostic agreement, stability checks and limitations together.

\section{System Model and Calibrated Surrogate Framework}

The framework evaluates coherent-link performance metrics under controlled PCS and baseline configurations. The current implementation should be interpreted as a calibrated surrogate workflow. In particular, the waveform and SSFM-like components use normalized or effective scaling choices and should not be presented as a fully physical Manakov/NLSE solver.

\subsection{Physical interpretation}

The model uses effective propagation and impairment parameters to study relative trends. Dispersion, nonlinearity, noise accumulation and learned mitigation are therefore interpreted as components of a calibrated surrogate workflow. This interpretation is different from a fully parameterized physical propagation model using independently specified $\beta_2$, $\gamma$, attenuation, amplifier gain, noise figure and Raman/SRS dynamics.

This distinction affects the validation language. Raw absolute GSNR should not be interpreted as directly validated against GNPy. Instead, GNPy is used as an external calibration diagnostic. The raw C-band comparison shows a conservative offset, while the offset-calibrated diagnostic tests whether the span-dependent trend is consistent after one transparent correction.

The manuscript therefore distinguishes:
\begin{itemize}
    \item trend-level PCS evaluation;
    \item raw absolute GSNR comparison;
    \item offset-calibrated external diagnostic agreement;
    \item limitations of the normalized surrogate representation.
\end{itemize}

\section{Baseline and Comparison Design}

The comparison design separates shaping, learned mitigation and external calibration effects. The minimum baseline set is:
\begin{itemize}
    \item \textbf{Uniform modulation:} the unshaped reference case.
    \item \textbf{PCS-only:} isolates shaping gain from learned mitigation.
    \item \textbf{Neural-only mitigation:} isolates receiver-side or mitigation gain without shaping, where available.
    \item \textbf{PCS plus neural mitigation:} evaluates the combined configuration, where available.
    \item \textbf{Raw GNPy reference:} reports the uncalibrated external mismatch.
    \item \textbf{Offset-calibrated GNPy diagnostic:} reports calibrated trend agreement after applying one transparent offset.
    \item \textbf{Computational cost:} reports runtime or complexity where available, so accuracy and cost are not mixed.
\end{itemize}

This comparison structure is designed to avoid attributing all gains to a single mechanism. It also ensures that calibration is reported explicitly rather than hidden inside the final metric.

\section{Validation Methodology}

The validation methodology contains three levels.

\subsection{Repeated-seed stability}

Repeated-seed evaluation is used to estimate whether PCS gains remain stable under stochastic simulation variation.

\subsection{Symbol-count stability}

Larger symbol-count checks are used to reduce the risk that PCS gains are artifacts of short simulated sequences.

\subsection{External GNPy calibration diagnostic}

External GNPy comparison is used as a calibration diagnostic, not as raw absolute validation. The raw C-band reference comparison showed a conservative offset between the proposed workflow and GNPy. A multi-span sweep then showed that this offset is mostly constant across span count. Applying a single offset yielded sub-dB diagnostic agreement.

\section{Results}

\subsection{C-band stronger stability result}

The Day-10 C-band stronger stability test used 65,536 symbols and 10 seeds. The high-count GMI gain was 0.0209167 bit/symbol, with an absolute GMI-gain drift of 0.0002379 and relative drift of 0.011375.

\subsection{Raw GNPy comparison}

The Day-11 raw C-band GNPy reference comparison did not pass as direct absolute validation. The GNPy reference GSNR was 18.780000 dB, while the O-E-S-C-L uniform GSNR was 12.863055 dB, producing a raw error of -5.916945 dB.

\subsection{Multi-span offset analysis}

The Day-12 multi-span alignment sweep found a mean O-E-S-C-L minus GNPy offset of -6.152843 dB, offset standard deviation of 0.328463 dB and error-versus-span slope of 0.104764 dB/span. This supports the interpretation of a mostly constant conservative offset.

\subsection{Offset-calibrated diagnostic}

The Day-13 offset-calibrated diagnostic applied a constant +6.152843 dB correction. The calibrated RMSE was 0.284458 dB and the leave-one-out RMSE was 0.379277 dB. These results support calibrated trend agreement but not raw absolute validation.

\begin{table}
\caption{Summary of calibrated external diagnostic validation. These results support calibrated trend agreement, not raw absolute GNPy validation.}
\label{tab:validation-summary}
\begin{tabular}{lllll}
\toprule
Validation item & RMSE (dB) & MAE (dB) & Max abs. error (dB) & Interpretation \\
\midrule
Day-13 offset-calibrated diagnostic & 0.284458 & 0.240163 & 0.480325 & Sub-dB calibrated diagnostic agreement \\
Day-13 leave-one-out diagnostic & 0.379277 & 0.320217 & 0.640434 & Sub-dB LOO calibrated diagnostic agreement \\
\bottomrule
\end{tabular}
\end{table}


\section{Discussion}

The results support the use of the framework for calibrated PCS trend analysis. They do not support claims of raw absolute GNPy validation or experimental validation. The main value of the workflow is transparent claim-controlled evaluation, where calibrated agreement and raw disagreement are both reported.

\section{Limitations}

The current framework has several limitations:
\begin{itemize}
    \item it is a calibrated surrogate rather than a fully physical Manakov/NLSE simulator;
    \item raw absolute GSNR agreement with GNPy is weak before calibration;
    \item broader independent validation across span count, baud rate, launch power, channel loading and fiber parameters is still needed;
    \item no laboratory or field measurement validation is claimed;
    \item calibrated agreement should not be interpreted as direct physical proof.
\end{itemize}

\section{Reproducibility}

The repository includes a root README, dependency lock file, smoke test, claim-control checker and curated validation reports. The smoke test checks the presence of key validation artifacts and verifies the Day-13 calibrated and leave-one-out RMSE values.

The manuscript tables are generated using:
\begin{verbatim}
python scripts\generate_manuscript_artifacts.py
\end{verbatim}

Repository hygiene and claim-control basics are checked using:
\begin{verbatim}
python scripts\smoke_test_reproducibility.py
python scripts\check_pnc_submission_readiness.py
\end{verbatim}

These checks do not imply that the manuscript is ready for submission; they only verify that the current repository state preserves the curated evidence and avoids selected overclaiming patterns.

\section{Conclusion}

This paper presents a calibrated surrogate framework for PCS trend evaluation in coherent optical link scenarios. The current validation evidence supports repeated-seed and symbol-count stability for selected PCS trends and shows that the C-band raw GNPy mismatch is largely explained by a near-constant conservative offset. After applying a single offset, the C-band multi-span trend aligns with GNPy with sub-dB error. The work should therefore be interpreted as a reproducibility-oriented calibrated surrogate study rather than a direct physical simulator or experimental validation study.

\section*{Data and Code Availability}

The code and curated validation artifacts will be made available in a public GitHub repository. A versioned archival DOI will be created before final submission.

\section*{Conflict of Interest}

The authors declare no conflict of interest.

\bibliography{references}

\end{document}
